% ══════════════════════════════════════════════════════════════
%  §6  Formal Properties and Discussion
% ══════════════════════════════════════════════════════════════
\section{Formal Properties and Discussion}
\label{sec:properties}

This section establishes the formal relationship between \textsc{Cir}
executions and \textsc{Cvn} firings, states the guarantees provided
by the analysis, and discusses the modeling decisions and limitations
of the approach.

% ─────────────────────────────────────────────────────────────
\subsection{CIR Execution Semantics}
\label{subsec:cir-semantics}

\begin{definition}[\textsc{Cir} configuration]
\label{def:cir-config}
A \textsc{Cir} configuration is a tuple
$C = (\mathit{pc}, L, W, \mathit{NW}, \mathit{NA}, V)$
where
\begin{itemize}[leftmargin=1.4em]
  \item $\mathit{pc} : \mathit{ThreadId} \to \mathit{SID} \cup
        \{\mathit{ret}\}$ maps each thread to its current statement,
  \item $L : \mathit{ResourceName} \to \mathbb{N}_0$ records the
        available count of each synchronization resource (mutexes,
        semaphores, channels) and each condvar notification token,
  \item $W : \mathit{SID} \to \mathbb{N}_0$ records the number of
        threads blocked at each wait site,
  \item $\mathit{NW} : \mathit{CondvarName} \to \mathbb{N}_0$
        records the total number of waiters for each condvar,
  \item $\mathit{NA} : \mathit{SID} \to \mathsf{Bool}$ records the
        notify-all flag for each wait site, and
  \item $V : \mathit{VarName} \to \mathit{Val}$ is the variable
        store.
\end{itemize}
\end{definition}

The \textsc{Cir} step rules for condvar operations are:

\textbf{Wait.} Thread $\theta$ at $(\mathit{sid},\;
\textsf{wait}(\mathit{cv}, m),\; \textsf{next}(\mathit{sid}'))$
performs two sub-steps. Enter: $L'(m) = L(m) + 1$,
$W'(\mathit{sid}) = W(\mathit{sid}) + 1$,
$\mathit{NW}'(\mathit{cv}) = \mathit{NW}(\mathit{cv}) + 1$,
$\mathit{NA}'(\mathit{sid}) = \mathsf{false}$,
thread enters waiting state at $\mathit{sid}$. Wake (by notify):
requires $L(\mathit{cv}) \geq 1$; sets $L'(\mathit{cv}) = L(\mathit{cv}) - 1$,
$W'(\mathit{sid}) = W(\mathit{sid}) - 1$,
$\mathit{NW}'(\mathit{cv}) = \mathit{NW}(\mathit{cv}) - 1$,
thread enters reacquire state. Wake (by notify-all):
requires $\mathit{NA}(\mathit{sid}) = \mathsf{true}$; sets
$W'(\mathit{sid}) = W(\mathit{sid}) - 1$,
$\mathit{NW}'(\mathit{cv}) = \mathit{NW}(\mathit{cv}) - 1$,
$\mathit{NA}'(\mathit{sid}) = \mathsf{false}$,
thread enters reacquire state. Reacquire: requires $L(m) \geq 1$;
sets $L'(m) = L(m) - 1$, $\mathit{pc}'(\theta) = \mathit{sid}'$.

\textbf{Notify\_one.} If $\mathit{NW}(\mathit{cv}) > 0$: set
$L'(\mathit{cv}) = L(\mathit{cv}) + 1$,
$\mathit{pc}'(\theta) = \mathit{sid}'$. If
$\mathit{NW}(\mathit{cv}) = 0$: notification lost,
$\mathit{pc}'(\theta) = \mathit{sid}'$.

\textbf{Notify\_all.} If $\mathit{NW}(\mathit{cv}) > 0$: for each
wait site $w_i$ of $\mathit{cv}$, set
$\mathit{NA}'(w_i) = \mathsf{true}$,
$\mathit{pc}'(\theta) = \mathit{sid}'$. If
$\mathit{NW}(\mathit{cv}) = 0$: notification lost,
$\mathit{pc}'(\theta) = \mathit{sid}'$.


% ─────────────────────────────────────────────────────────────
\subsection{Encoding Function}
\label{subsec:encoding}

\begin{definition}[Encoding]
\label{def:encoding}
The encoding function $\mathsf{enc}$ maps a \textsc{Cir}
configuration $C = (\mathit{pc}, L, W, \mathit{NW}, \mathit{NA}, V)$
to a \textsc{Cvn} state $(M, V')$ as follows:
\begin{itemize}[leftmargin=1.4em]
  \item For each thread $\theta$ at $\mathit{sid}$ in function $f$
        (not waiting, not reacquiring):
        $M(\mathit{cp}(f, \mathit{sid})) \mathrel{+}= 1$.
  \item For each thread $\theta$ waiting at site $\mathit{sid}$:
        $M(\mathit{wp}(\mathit{sid})) \mathrel{+}= 1$.
  \item For each thread $\theta$ reacquiring at site $\mathit{sid}$:
        $M(\mathit{ra}(\mathit{sid})) \mathrel{+}= 1$.
  \item For each synchronization resource $r$:
        $M(\mathit{rp}(r)) = L(r)$.
  \item For each condvar $\mathit{cv}$:
        $M(\mathit{rp}(\mathit{cv})) = L(\mathit{cv})$.
  \item $V'$ extends $V$ with:
        $V'[\mathit{nw}_{\mathit{cv}}] = \mathit{NW}(\mathit{cv})$
        for each condvar, and
        $V'[\mathit{na}_{w_i}] = \mathit{NA}(w_i)$ for each wait
        site.
\end{itemize}
\end{definition}

The encoding preserves all concurrency-relevant information. It does
not preserve thread identity (which thread is at which position),
because \textsc{Cvn} tokens are indistinguishable. This is
acceptable because the analysis does not require thread identity:
two tokens in the same control place represent two interchangeable
threads at the same program point.

% ─────────────────────────────────────────────────────────────

\subsection{Step Correspondence}
\label{subsec:step-correspondence}

\begin{theorem}[Forward simulation]
\label{thm:forward}
Let $C$ be a \textsc{Cir} configuration and
$(M, V) = \mathsf{enc}(C)$. For every \textsc{Cir} step
$C \xrightarrow{o} C'$, there exists a \textsc{Cvn} transition $t$
such that $(M, V) \xrightarrow{t} (M', V')$ and
$(M', V') = \mathsf{enc}(C')$.
\end{theorem}

\begin{proof}
By case analysis on the operation $o$. We present the condvar cases;
the remaining cases follow the same structure as before.

\textbf{Case} $o = \textsf{lock}(m)$. The \textsc{Cir} step
requires $L(m) \geq 1$ and produces $L'(m) = L(m) - 1$,
$\mathit{pc}'(\theta) = \mathit{sid}'$. The corresponding
\textsc{Cvn} transition $t$ [Lock] has input arcs from
$\mathit{cp}(f, \mathit{sid})$ weight 1 and $\mathit{rp}(m)$
weight 1, and output arc to $\mathit{cp}(f, \mathit{sid}')$
weight 1. The precondition $M(\mathit{rp}(m)) = L(m) \geq 1$
guarantees that $t$ is enabled. Firing $t$ produces
$M'(\mathit{rp}(m)) = M(\mathit{rp}(m)) - 1 = L'(m)$ and
$M'(\mathit{cp}(f, \mathit{sid}')) = M(\mathit{cp}(f,
\mathit{sid}')) + 1$, matching $\mathsf{enc}(C')$.

\textbf{Case} $o = \textsf{drop}(m)$. Symmetric: the transition
returns a token to $\mathit{rp}(m)$.

\textbf{Case} $o = \textsf{write}(x, v)$. The \textsc{Cir} step
sets $V'[x] = v$. The \textsc{Cvn} transition [VarWrite] carries
the update $V[x] \leftarrow v$, producing the same variable
store.

\textbf{Case} $o = \textsf{read}(x)$ with
$\textsf{branch}(\beta, \mathit{sid}_t, \mathit{sid}_f)$. The
\textsc{Cir} step evaluates $\beta$ to a definite boolean and
branches accordingly. The \textsc{Cvn} has two transitions
$t_\top$ and $t_\bot$ with complementary guards. Since $\beta$
evaluates to a concrete boolean, exactly one guard evaluates to
$\mathsf{true}$ and the other to $\mathsf{false}$. The enabled
transition matches the \textsc{Cir} branch direction.

\textbf{Case} $o = \textsf{wait}(\mathit{cv}, m)$. The
\textsc{Cir} step moves the thread to a waiting state and releases
the mutex. The \textsc{Cvn} transition $t_w$ [CondvarWait] moves
the token from $\mathit{cp}(f, \mathit{sid})$ to
$\mathit{wp}(\mathit{sid})$ and returns a token to
$\mathit{rp}(m)$, matching the encoding of the new configuration.

\textbf{Case} $o = \textsf{notify\_one}(\mathit{cv})$. If at
least one thread is waiting (some $W(\mathit{sid}_w) \geq 1$),
the \textsc{Cir} step wakes one waiter nondeterministically. The
corresponding wake transition $t_i$ [CondvarNotify] consumes a
token from $\mathit{wp}(\mathit{sid}_w)$, matching the waiter
selection. If no thread is waiting, the \textsc{Cir} step has no
wake effect and advances control. The corresponding pass-through
transition advances control without consuming from any wait place.

\textbf{Case} $o = \textsf{spawn}(g)$. The \textsc{Cir} step
creates a new thread at the entry of $g$. The \textsc{Cvn}
transition [Spawn] produces tokens in both
$\mathit{cp}(f, \mathit{sid}')$ and
$\mathit{cp}(g, \mathit{sid}_0^g)$, matching the fork.

\textbf{Case} $o = \textsf{wait}(\mathit{cv}, m)$, \emph{enter
sub-step}. The \textsc{Cir} step increments
$W(\mathit{sid})$, increments $\mathit{NW}(\mathit{cv})$,
resets $\mathit{NA}(\mathit{sid})$, and releases the mutex
($L(m) + 1$). The \textsc{Cvn} transition
$t_{\mathit{enter}}$ [CondvarWaitEnter] consumes one token from
$\mathit{cp}(f, \mathit{sid})$, produces one token in
$\mathit{wp}(\mathit{sid})$ and one token in $\mathit{rp}(m)$,
updates $V[\mathit{nw}_{\mathit{cv}}] \leftarrow
V[\mathit{nw}_{\mathit{cv}}] + 1$ and
$V[\mathit{na}_{\mathit{sid}}] \leftarrow \mathsf{false}$.
Enabling: $M(\mathit{cp}(f, \mathit{sid})) \geq 1$ holds because
the thread is at $\mathit{sid}$. No guard. The resulting marking and
variable store match $\mathsf{enc}(C')$.

\textbf{Case} $o = \textsf{wait}(\mathit{cv}, m)$, \emph{wake by
notify sub-step}. The \textsc{Cir} step requires
$L(\mathit{cv}) \geq 1$, decrements $W(\mathit{sid})$, decrements
$\mathit{NW}(\mathit{cv})$, and consumes one notification token. The
\textsc{Cvn} transition $t_{\mathit{wake1}}$
[CondvarWakeByNotify] consumes one token from
$\mathit{wp}(\mathit{sid})$ and one token from
$\mathit{rp}(\mathit{cv})$, produces one token in
$\mathit{ra}(\mathit{sid})$, updates
$V[\mathit{nw}_{\mathit{cv}}] \leftarrow
V[\mathit{nw}_{\mathit{cv}}] - 1$. Enabling:
$M(\mathit{wp}(\mathit{sid})) \geq 1$ holds because the thread is
waiting, and $M(\mathit{rp}(\mathit{cv})) = L(\mathit{cv}) \geq 1$.
No guard. Match follows.

\textbf{Case} $o = \textsf{wait}(\mathit{cv}, m)$, \emph{wake by
notify-all sub-step}. The \textsc{Cir} step requires
$\mathit{NA}(\mathit{sid}) = \mathsf{true}$, does not consume a
notification token, decrements $W(\mathit{sid})$ and
$\mathit{NW}(\mathit{cv})$, resets
$\mathit{NA}(\mathit{sid})$. The \textsc{Cvn} transition
$t_{\mathit{wakeA}}$ [CondvarWakeByNotifyAll] consumes one token
from $\mathit{wp}(\mathit{sid})$, has guard
$V[\mathit{na}_{\mathit{sid}}] = \mathsf{true}$, produces one
token in $\mathit{ra}(\mathit{sid})$, and updates
$V[\mathit{nw}_{\mathit{cv}}] \leftarrow
V[\mathit{nw}_{\mathit{cv}}] - 1$,
$V[\mathit{na}_{\mathit{sid}}] \leftarrow \mathsf{false}$. Enabling:
token in $\mathit{wp}(\mathit{sid})$ holds; guard holds because
$V[\mathit{na}_{\mathit{sid}}] = \mathit{NA}(\mathit{sid}) =
\mathsf{true}$. Match follows.

\textbf{Case} $o = \textsf{wait}(\mathit{cv}, m)$, \emph{reacquire
sub-step}. The \textsc{Cir} step requires $L(m) \geq 1$ and sets
$\mathit{pc}'(\theta) = \mathit{sid}'$. The \textsc{Cvn} transition
$t_{\mathit{reacq}}$ [CondvarReacquire] consumes one token from
$\mathit{ra}(\mathit{sid})$ and one from $\mathit{rp}(m)$, produces
one token in $\mathit{cp}(f, \mathit{sid}')$. Enabling: both tokens
are present. Match follows.

\textbf{Case} $o = \textsf{notify\_one}(\mathit{cv})$ with
$\mathit{NW}(\mathit{cv}) > 0$. The \textsc{Cir} step deposits one
notification token: $L'(\mathit{cv}) = L(\mathit{cv}) + 1$. The
\textsc{Cvn} transition $t_{\mathit{notify}}$ [CondvarNotify] has
guard $V[\mathit{nw}_{\mathit{cv}}] > 0$, which holds because
$V[\mathit{nw}_{\mathit{cv}}] = \mathit{NW}(\mathit{cv}) > 0$.
The transition produces one token in $\mathit{rp}(\mathit{cv})$.
Match follows.

\textbf{Case} $o = \textsf{notify\_one}(\mathit{cv})$ with
$\mathit{NW}(\mathit{cv}) = 0$. The notification is lost. The
\textsc{Cvn} transition $t_{\mathit{lost}}$ [CondvarNotifyLost]
has guard $V[\mathit{nw}_{\mathit{cv}}] = 0$, which holds. The
transition advances control without depositing a token. Match
follows.

\textbf{Case} $o = \textsf{notify\_all}(\mathit{cv})$ with
$\mathit{NW}(\mathit{cv}) > 0$. The \textsc{Cir} step sets
$\mathit{NA}'(w_i) = \mathsf{true}$ for all wait sites $w_i$ of
$\mathit{cv}$. The \textsc{Cvn} transition
$t_{\mathit{notifyAll}}$ [CondvarNotifyAll] has guard
$V[\mathit{nw}_{\mathit{cv}}] > 0$, which holds. The transition
updates $V[\mathit{na}_{w_i}] \leftarrow \mathsf{true}$ for all
$w_i$. Match follows.

\textbf{Case} $o = \textsf{notify\_all}(\mathit{cv})$ with
$\mathit{NW}(\mathit{cv}) = 0$. Same as the notify-lost case.

The remaining cases (\textsf{join}, \textsf{cas}, \textsf{acquire},
\textsf{release}, \textsf{send}, \textsf{recv}, \textsf{call},
\textsf{return}) follow the same pattern. Each \textsc{Cir}
precondition maps to the corresponding \textsc{Cvn} enabling
condition, and each postcondition maps to the corresponding firing
effect.
\end{proof}

The forward simulation guarantees that every \textsc{Cir} behavior
is represented in the \textsc{Cvn}. The \textsc{Cvn} cannot miss a
bug that exists in the \textsc{Cir}.

\begin{theorem}[Backward simulation under concrete valuation]
\label{thm:backward}
Let $(M, V) = \mathsf{enc}(C)$ where all variables in $V$ have
concrete values (no variable equals $\top$). For every enabled
\textsc{Cvn} transition $t$ at $(M, V)$, there exists a \textsc{Cir}
step $C \xrightarrow{o} C'$ such that
$\mathsf{enc}(C') = \mathsf{fire}(t, M, V)$.
\end{theorem}

\begin{proof}
When all variables are concrete, guard evaluation in the \textsc{Cvn}
produces only $\mathsf{true}$ or $\mathsf{false}$, never
$\mathsf{unknown}$. The enabling conditions of \textsc{Cvn}
transitions therefore coincide exactly with the preconditions of
\textsc{Cir} operations. The argument proceeds by case analysis,
symmetric to Theorem~\ref{thm:forward}.
\end{proof}

When some variable has the value $\top$, the backward simulation may
fail: a guard involving that variable evaluates to $\mathsf{unknown}$,
enabling the transition in the \textsc{Cvn}, but the corresponding
\textsc{Cir} branch may not be takeable under any concrete valuation.
This is the source of potential false positives, discussed in
Section~\ref{subsec:over-approx}.
% ─────────────────────────────────────────────────────────────
\subsection{Deadlock Correspondence}
\label{subsec:deadlock-correspondence}

\begin{theorem}[Deadlock correspondence]
\label{thm:deadlock-corr}
A \textsc{Cir} configuration $C$ is a deadlock (non-terminal with no
enabled step) if and only if $\mathsf{enc}(C)$ is a deadlock in the
\textsc{Cvn} (non-terminal marking with no enabled transition),
provided all variables have concrete values.
\end{theorem}

\begin{proof}
A configuration is a deadlock when the set of enabled
operations/transitions is empty and at least one thread has not
returned. By Theorems~\ref{thm:forward} and~\ref{thm:backward}, the
enabled sets coincide under concrete valuation. The non-terminal
condition is preserved by the encoding: a thread at
$\mathit{sid} \neq \mathit{ret}$ maps to a token in a non-return
control place.
\end{proof}

Under $\top$-valued variables, the correspondence is one-directional:
a \textsc{Cir} deadlock implies a \textsc{Cvn} deadlock (by
Theorem~\ref{thm:forward}), but a \textsc{Cvn} deadlock may not
correspond to any \textsc{Cir} deadlock (the \textsc{Cvn} may have
reached the state through an infeasible branch enabled by
$\mathsf{unknown}$).


\subsection{Signal Loss Soundness}
\label{subsec:signal-loss-soundness}

The updated condvar model enables a precise characterization of
signal loss.

\begin{proposition}[Signal loss detection]
\label{prop:signal-loss}
A \textsf{notify\_one} or \textsf{notify\_all} operation at
$\mathit{sid}_n$ causes a signal loss in a \textsc{Cir} execution if
and only if the corresponding \textsc{Cvn} transition
$t_{\mathit{lost}}$ or $t_{\mathit{allLost}}$ fires at state
$(M, V)$ with $V[\mathit{nw}_{\mathit{cv}}] = 0$.
\end{proposition}

\begin{proof}
Forward direction: if the \textsc{Cir} notification occurs with
$\mathit{NW}(\mathit{cv}) = 0$, then by the encoding
$V[\mathit{nw}_{\mathit{cv}}] = 0$, and the guard of
$t_{\mathit{lost}}$ is satisfied. Backward direction: if
$t_{\mathit{lost}}$ fires, then
$V[\mathit{nw}_{\mathit{cv}}] = 0 = \mathit{NW}(\mathit{cv})$,
and the \textsc{Cir} notification has no waiters.
\end{proof}

This proposition establishes that the signal-loss annotation on the
lost transitions is both sound (no false negatives) and precise (no
false positives) with respect to the \textsc{Cir} semantics. A
signal loss is recorded if and only if the notification has no
waiters in the corresponding \textsc{Cir} execution.



% ─────────────────────────────────────────────────────────────
\subsection{Soundness of Over-Approximation}
\label{subsec:over-approx}

\begin{theorem}[Soundness]
\label{thm:soundness}
If a \textsc{Cir} execution can reach a bug configuration (deadlock,
signal loss, starvation, or livelock), then the \textsc{Cvn}
state-space search will find a corresponding bug state.
\end{theorem}

\begin{proof}
Let $C_0 \xrightarrow{o_1} C_1 \xrightarrow{o_2} \cdots
\xrightarrow{o_n} C_n$ be a \textsc{Cir} execution reaching a bug
configuration $C_n$. By Theorem~\ref{thm:forward}, there exists a
corresponding \textsc{Cvn} firing sequence
$(M_0, V_0) \xrightarrow{t_1} \cdots \xrightarrow{t_n}
(M_n, V_n)$ with $(M_i, V_i) = \mathsf{enc}(C_i)$ for each $i$.

The $\mathsf{unknown}$ mechanism only adds enabled transitions to
the \textsc{Cvn} (both branches of an unknown guard become enabled).
It never removes transitions that would be enabled under concrete
evaluation. Therefore, every transition enabled in the concrete
\textsc{Cvn} is also enabled in the $\mathsf{unknown}$-augmented
\textsc{Cvn}. The firing sequence above remains valid.

The bug predicate (deadlock, signal loss, starvation, livelock) is
defined over \textsc{Cvn} states and state-graph structure
(Definitions~\ref{def:deadlock}--\ref{def:channel-block}). Since
$(M_n, V_n) = \mathsf{enc}(C_n)$ and $C_n$ satisfies the bug
predicate at the \textsc{Cir} level, $(M_n, V_n)$ satisfies the
corresponding predicate at the \textsc{Cvn} level.

The finite state space (Theorem~\ref{thm:finite}) guarantees that
the search terminates and visits $(M_n, V_n)$.
\end{proof}

The converse does not hold: the \textsc{Cvn} may report bugs in
states that are reachable only through $\mathsf{unknown}$-enabled
branches and that no concrete \textsc{Cir} execution can reach.
These are false positives. In practice, false positives arise only
when variables modified by summarized functions
(Section~\ref{subsec:translation}, Phase~3) are used in branch
conditions. For fully modeled programs where all functions have
bodies in the \textsc{Cir}, all variables have concrete values
throughout execution, and the analysis produces no false positives.

% ─────────────────────────────────────────────────────────────
\subsection{Loop Modeling}
\label{subsec:loop-discussion}

The system does not perform automated loop unrolling. This is a
deliberate design decision based on two observations.

\textbf{Concurrency bugs are topological.} Deadlocks, signal losses,
and starvation are triggered by the structure of synchronization
interactions, not by loop iteration counts. A circular lock
dependency causes deadlock on the first iteration. A misplaced
notification causes signal loss regardless of how many times the loop
body executes. Two or three concurrent instances are sufficient to
expose the vast majority of synchronization bugs.

\textbf{LLM-driven instance unfolding replaces loop unrolling.}
When the source code contains a loop that spawns $N$ workers, the
LLM unfolds the loop into a small number (typically 2 or 3) of
explicit \textsf{spawn} statements in the \textsc{Cir}. Each
instance receives a unique function name (e.g.,
\texttt{worker\_0}, \texttt{worker\_1}) and therefore a unique set
of control places in the \textsc{Cvn}. Black tokens suffice to
distinguish thread positions because the places themselves are
distinct.

Non-spawn loops retain their back-edge structure in \textsc{Cir}. A
loop body that modifies a shared variable may cause that variable to
take different values on different iterations. If the variable's
value cannot be determined statically, it becomes $\top$ in the
\textsc{Cvn} after the first iteration. The three-valued guard
evaluation then enables both branches of any condition involving
that variable, ensuring that all paths through the loop body are
explored. This provides a safe over-approximation of all loop
iterations without explicit unrolling.

\textbf{Consequence.} The system verifies a fixed, bounded number of
concurrent instances, not parametric correctness for arbitrary $N$.
Bugs that require more instances than the LLM generates (e.g., a
$k$-way circular deadlock where $k$ exceeds the unfolding count)
will not be detected. In practice, setting the unfolding count to
$\min(N, 3)$ is sufficient for the patterns in our test matrix
(Section~\ref{sec:evaluation}).

% ─────────────────────────────────────────────────────────────
\subsection{Why Colored Tokens Are Unnecessary}
\label{subsec:no-color}

Classical colored Petri nets (CPNs) attach typed data to tokens,
enabling a single parameterized net to model multiple instances. In
our architecture, this parameterization is unnecessary for two
reasons.

First, the LLM has already resolved all resource identities at the
\textsc{Cir} level. Tokens do not need to carry resource names
because each resource has a dedicated place. Second, the LLM has
already unfolded concurrent instances into distinct functions. Tokens
do not need to carry thread identities because each instance has
dedicated control places. Black tokens are sufficient.

This simplification eliminates the need for color sets, binding
expressions, arc inscriptions, and the canonical state
representation that CPNs require for state-space exploration. The
resulting \textsc{Cvn} is a standard weighted P/T net (augmented
with global-variable guards), for which efficient analysis algorithms
are well established~\cite{Murata1989}.

% ─────────────────────────────────────────────────────────────
\subsection{Limitations}
\label{subsec:limitations}

\textbf{LLM reliability.} The correctness of the \textsc{Cir}
depends on the LLM's understanding of the source code. Complex
aliasing through unsafe code, raw pointers, or foreign function
interfaces may exceed current LLM capabilities. Errors in the LLM's
resource identification manifest as structural errors in \textsc{Cir}
and are caught by the static checker, but subtle omissions (e.g.,
failing to model a shared resource) may lead to missed bugs.

\textbf{Over-approximation.} The $\top$ mechanism may produce false
positives when branch conditions involve variables whose values
cannot be determined. The false positive rate depends on the
proportion of summarized functions and the complexity of branch
conditions.

\textbf{Bounded instances.} The system verifies a fixed number of
unfolded instances. Parametric correctness (``safe for any $N$'') is
outside the current scope.

\textbf{No weak-memory model.} Atomic operations are modeled at the
variable level without a formal memory model. Bugs that depend on
weak-memory ordering (e.g., missing barriers on ARM) are not
detected.

\textbf{State explosion.} The state space grows exponentially with
the number of concurrent entities. For programs with many threads
and many synchronization operations, the search may exceed the
configured state limit before completing. The timeout and
state-limit parameters allow trading completeness for practicality.

